% ============================================================
%  Vibration Analysis Suite — User Manual
%  Author : Akshat_EA
%  Engine : pdfLaTeX  (compile twice for TOC + refs)
% ============================================================
\documentclass[12pt, a4paper]{article}

% ── Packages ────────────────────────────────────────────────
\usepackage[a4paper, top=2.5cm, bottom=2.5cm,
            left=3cm,  right=3cm]{geometry}
\usepackage{fontenc}
\usepackage{inputenc}
\usepackage{lmodern}
\usepackage{microtype}
\usepackage{parskip}
\usepackage{titlesec}
\usepackage{titletoc}
\usepackage{graphicx}
\usepackage{xcolor}
\usepackage{mdframed}
\usepackage{enumitem}
\usepackage{booktabs}
\usepackage{array}
\usepackage{tabularx}
\usepackage{fancyhdr}
\usepackage{hyperref}
\usepackage{listings}
\usepackage{tcolorbox}
\usepackage{amsmath}
\usepackage{amssymb}
\usepackage{float}
\usepackage{caption}
\usepackage{setspace}
\usepackage{eso-pic}

% ── Colour palette (Catppuccin-inspired) ────────────────────
\definecolor{accent}    {HTML}{1e66f5}
\definecolor{accentdark}{HTML}{0a50e0}
\definecolor{muted}     {HTML}{6c7086}
\definecolor{warn}      {HTML}{d20f39}
\definecolor{success}   {HTML}{40a02b}
\definecolor{codebg}    {HTML}{f4f4f8}
\definecolor{codefg}    {HTML}{1e1e2e}
\definecolor{notebg}    {HTML}{dce0f4}
\definecolor{warnbg}    {HTML}{fce8ec}
\definecolor{tipbg}     {HTML}{e6f4ea}
\definecolor{headerbg}  {HTML}{1e1e2e}
\definecolor{headertext}{HTML}{cdd6f4}

% ── Hyperref setup ───────────────────────────────────────────
\hypersetup{
  colorlinks = true,
  linkcolor  = accent,
  urlcolor   = accentdark,
  citecolor  = muted,
  pdftitle   = {Vibration Analysis Suite — User Manual},
  pdfauthor  = {Akshat\_EA},
}

% ── Section title styling ────────────────────────────────────
\titleformat{\section}
  {\Large\bfseries\color{accent}}
  {\thesection}{1em}{}
  [\color{accent}\titlerule]

\titleformat{\subsection}
  {\large\bfseries\color{accentdark}}
  {\thesubsection}{1em}{}

\titleformat{\subsubsection}
  {\normalsize\bfseries\color{codefg}}
  {\thesubsubsection}{1em}{}

% ── Header / Footer ─────────────────────────────────────────
\pagestyle{fancy}
\fancyhf{}
\fancyhead[L]{\small\color{muted}\textit{Vibration Analysis Suite}}
\fancyhead[R]{\small\color{muted}\textit{User Manual}}
\fancyfoot[C]{\small\color{muted}\thepage}
\renewcommand{\headrulewidth}{0.4pt}
\renewcommand{\footrulewidth}{0pt}

% ── Code listing style ───────────────────────────────────────
\lstdefinestyle{shell}{
  backgroundcolor = \color{codebg},
  basicstyle      = \ttfamily\small\color{codefg},
  breaklines      = true,
  frame           = single,
  framerule       = 0.6pt,
  rulecolor       = \color{accent},
  xleftmargin     = 8pt,
  xrightmargin    = 8pt,
  aboveskip       = 8pt,
  belowskip       = 4pt,
}
\lstset{style=shell}

% ── Custom boxes ─────────────────────────────────────────────
\tcbuselibrary{skins, breakable}

\newtcolorbox{notebox}{
  enhanced, breakable,
  colback  = notebg,
  colframe = accent,
  fonttitle= \bfseries\color{accent},
  title    = {$\blacktriangleright$\; Note},
  boxrule  = 0.8pt,
  arc      = 4pt,
  left=8pt, right=8pt, top=6pt, bottom=6pt,
}

\newtcolorbox{warnbox}{
  enhanced, breakable,
  colback  = warnbg,
  colframe = warn,
  fonttitle= \bfseries\color{warn},
  title    = {$\bigtriangleup$\; Warning},
  boxrule  = 0.8pt,
  arc      = 4pt,
  left=8pt, right=8pt, top=6pt, bottom=6pt,
}

\newtcolorbox{tipbox}{
  enhanced, breakable,
  colback  = tipbg,
  colframe = success,
  fonttitle= \bfseries\color{success},
  title    = {$\checkmark$\; Tip},
  boxrule  = 0.8pt,
  arc      = 4pt,
  left=8pt, right=8pt, top=6pt, bottom=6pt,
}

% ── Step counter ─────────────────────────────────────────────
\newcounter{stepcnt}
\newcommand{\step}[1]{%
  \stepcounter{stepcnt}%
  \noindent\textbf{\color{accent}Step \thestepcnt\,:}\enspace #1\par\smallskip
}
\newcommand{\resetsteps}{\setcounter{stepcnt}{0}}

% ── Key shortcut macro ───────────────────────────────────────
\newcommand{\key}[1]{%
  \fbox{\small\ttfamily #1}%
}

% ── Button macro ─────────────────────────────────────────────
\newcommand{\btn}[1]{%
  \textbf{\color{accentdark}[#1]}%
}

% ============================================================
\begin{document}

% ── Cover Page ───────────────────────────────────────────────
\begin{titlepage}
\pagecolor{headerbg}
\color{headertext}
\vspace*{2cm}
\begin{center}
  {\fontsize{36}{44}\selectfont\bfseries Vibration Analysis Suite}\\[0.6cm]
  {\Large\color{muted} Complete User Manual}\\[2.5cm]

  \begin{mdframed}[linecolor=accent, linewidth=1.5pt,
                   backgroundcolor=headerbg,
                   leftmargin=2cm, rightmargin=2cm,
                   innertopmargin=12pt, innerbottommargin=12pt]
    \centering
    {\large Covering:}\\[0.4cm]
    {\normalsize
      Raw Crop \quad$\diamond$\quad FFT Analyzer \quad$\diamond$\quad Spectrogram
    }
  \end{mdframed}

  \vfill
  {\large\bfseries Developed by\;}{\large\color{accent}\texttt{Akshat\_EA}}\\[0.4cm]
  {\color{muted}\small Version 1.0 \quad$\cdot$\quad \today}
\end{center}
\end{titlepage}
\nopagecolor

% ── Table of Contents ────────────────────────────────────────
\newpage
\tableofcontents
\newpage

% ============================================================
\section{Introduction}
% ============================================================

The \textbf{Vibration Analysis Suite} is a desktop application developed by
\texttt{Akshat\_EA} for analysing tri-axial accelerometer data recorded in
binary \texttt{.bin} files.  It is the Python/PyQt6 replacement for three
MATLAB scripts previously used for the same purpose:

\begin{center}
\begin{tabular}{lll}
\toprule
\textbf{Original Script} & \textbf{Suite Tab} & \textbf{Purpose}\\
\midrule
\texttt{RawCrop.m}           & Raw Crop       & Load, inspect and crop recordings \\
\texttt{VibeDriver.m}        & FFT Analyzer   & Multi-file FFT and time-domain overlay \\
\texttt{Spectrogram\_SFFT.m} & Spectrogram    & White-hot spectrogram visualisation \\
\bottomrule
\end{tabular}
\end{center}

\subsection{Binary File Format}

Each \texttt{.bin} file stores accelerometer records in a fixed 18-byte
little-endian structure:

\begin{center}
\begin{tabular}{llll}
\toprule
\textbf{Field} & \textbf{Type} & \textbf{Bytes} & \textbf{Unit}\\
\midrule
\texttt{time\_ms} & \texttt{uint32}  & 4 & milliseconds \\
\texttt{accX}     & \texttt{float32} & 4 & milli-g (converted to g on load) \\
\texttt{accY}     & \texttt{float32} & 4 & milli-g (converted to g on load) \\
\texttt{accZ}     & \texttt{float32} & 4 & milli-g (converted to g on load) \\
\texttt{rate\_hz} & \texttt{uint16}  & 2 & Hz \\
\bottomrule
\end{tabular}
\end{center}

\begin{notebox}
Acceleration values are stored as milli-g in the file.  The application
automatically divides them by 1000 on load, so all displayed values are
in units of \textbf{g}.
\end{notebox}

% ============================================================
\section{Accessing the Machine}
% ============================================================

\subsection{Physical Access}

\begin{enumerate}[leftmargin=*, label=\textbf{\arabic*.}]
  \item Locate the designated analysis workstation.
  \item Ensure the machine is connected to power and the monitor is on.
  \item Press the \key{Power} button if the machine is off; if it is
        sleeping, press any key or move the mouse to wake it.
\end{enumerate}

\subsection{Booting into Windows}

\resetsteps
\step{When the machine powers on, watch for the manufacturer splash screen.}
\step{If a \textbf{Boot Menu} appears (common on shared or dual-boot machines),
      use the arrow keys \key{$\uparrow$} \key{$\downarrow$} to highlight
      \textbf{Windows 11} and press \key{Enter}.}
\step{Wait for the Windows login screen to appear.}
\step{Enter the Windows account credentials provided by \texttt{Akshat\_EA} and
      press \key{Enter} to log in.}
\step{Wait for the desktop to fully load before proceeding.}

\begin{warnbox}
Do \textbf{not} interrupt the boot process or force-shutdown the machine
while Windows is loading.  This can corrupt system files.
\end{warnbox}

% ============================================================
\section{Launching the Application}
% ============================================================

\subsection{Opening Windows PowerShell}

\resetsteps
\step{Right-click on the \textbf{Start} button (bottom-left of the taskbar).}
\step{Click \textbf{Terminal} or \textbf{Windows PowerShell} from the context
      menu.  A terminal window with a blue or dark background will appear.}

Alternatively:
\begin{itemize}
  \item Press \key{Win} + \key{R}, type \texttt{powershell}, and press \key{Enter}.
  \item Or press \key{Win} + \key{X} and select \textbf{Terminal}.
\end{itemize}

\subsection{Starting the Application}

Once the terminal is open, type the following command and press \key{Enter}:

\begin{lstlisting}
vibe
\end{lstlisting}

\begin{notebox}
The \texttt{vibe} command is a PowerShell alias pre-configured on this
machine by \texttt{Akshat\_EA}.  It internally runs:
\begin{lstlisting}
python "C:\Users\Akshat\Desktop\VibrationAnalysis\main.py"
\end{lstlisting}
You do \textbf{not} need to navigate to any folder or type the full path.
\end{notebox}

\begin{tipbox}
If the terminal shows \texttt{vibe : The term 'vibe' is not recognized},
contact \texttt{Akshat\_EA} to restore the PowerShell alias.
\end{tipbox}

% ============================================================
\section{Password Authentication}
% ============================================================

Immediately after launching, a password dialog appears before the main
application window.  The application \textbf{cannot} be used without
entering the correct password.

\subsection{Entering the Password}

\resetsteps
\step{The dialog box titled \textbf{Vibration Analysis Suite} will appear.}
\step{Click inside the password field (it may already be focused).}
\step{Type the password provided by \texttt{Akshat\_EA}.}
\step{Press \key{Enter} or click the \btn{Unlock} button.}
\step{If the password is correct, the main application window opens immediately.}

\subsection{Failed Attempts}

\begin{itemize}
  \item After each wrong password, the dialog shows how many attempts remain.
  \item You are allowed \textbf{3 attempts} in total.
  \item After 3 incorrect entries, the password field and the \btn{Unlock}
        button are disabled (greyed out), the message
        \textit{``Too many wrong attempts.  Contact Akshat\_EA for further
        information.''} is displayed, and the application closes
        automatically after a short delay.
\end{itemize}

\begin{warnbox}
If you exhaust all 3 attempts, you must contact \textbf{Akshat\_EA} to
regain access.  Re-launching with \texttt{vibe} immediately resets the
counter.
\end{warnbox}

% ============================================================
\section{Application Overview}
% ============================================================

After a successful login, the main window opens with the following layout:

\begin{center}
\begin{tabular}{ll}
\toprule
\textbf{Region} & \textbf{Description}\\
\midrule
Header bar    & App title, \textit{developed by Akshat\_EA} credit, theme toggle \\
Tab bar       & Three tabs: Raw Crop, FFT Analyzer, Spectrogram \\
Left panel    & Controls for the active tab (always 250--270 px wide) \\
Right panel   & Matplotlib plot canvas with navigation toolbar \\
Console       & Colour-coded log at the bottom of the window \\
\bottomrule
\end{tabular}
\end{center}

\subsection{Theme Toggle}

Click the \btn{$\odot$ Light Mode} / \btn{$\bullet$ Dark Mode} button in
the top-right of the header bar to switch between the light and dark
colour themes at any time.  The switch is \textbf{non-destructive}: it
does not reset zoom, pan, selections, or any loaded data.

\subsection{Console Log}

The console at the bottom of the window records every action taken.
Messages are colour-coded by tab:

\begin{itemize}
  \item \textcolor{accent}{\textbf{Blue}} — Raw Crop messages (\texttt{[Crop]})
  \item \textcolor{success}{\textbf{Green}} — FFT Analyzer messages (\texttt{[FFT]})
  \item \textcolor[HTML]{6a1b9a}{\textbf{Purple}} — Spectrogram messages (\texttt{[Spec]})
\end{itemize}

% ============================================================
\section{Raw Crop Tab}
% ============================================================

The Raw Crop tab is used to load a \texttt{.bin} file, visually inspect the
full time-domain signal on all three axes, select a sub-range of interest
using interactive trim bars, and export the cropped data as a new
\texttt{.bin} file.

\subsection{Loading a File}

\resetsteps
\step{Click \btn{Load BIN File}.}
\step{A file browser opens.  Navigate to the folder containing your
      \texttt{.bin} file and double-click it (or select and press \key{Open}).}
\step{Three time-domain subplots appear — one for each axis (AccX, AccY, AccZ).
      The left panel shows the file name, total record count, time range, and
      approximate sample rate.}

\subsection{Understanding the Plot}

Each subplot shows \textbf{acceleration (g)} on the Y-axis versus
\textbf{time (s)} on the X-axis.  Two coloured vertical dashed lines overlay
every subplot simultaneously:

\begin{itemize}
  \item \textcolor{success}{\textbf{Green dashed line}} — Start of the selected crop window
  \item \textcolor{warn}{\textbf{Red dashed line}} — End of the selected crop window
  \item A \textcolor{success}{semi-transparent green band} fills the selected region across all plots.
\end{itemize}

\subsection{Selecting a Crop Range}

There are two ways to set the start and end times:

\subsubsection{Method A — Drag the Trim Bars}

\begin{enumerate}[leftmargin=*, label=\textbf{\arabic*.}]
  \item Hover the mouse over either the green or red dashed line on any subplot.
        The cursor changes to a \textbf{horizontal resize arrow} ($\leftrightarrow$)
        when close enough to grab it (within 8 pixels).
  \item Click and hold the left mouse button, then drag left or right.
  \item All three subplots update simultaneously — the trim bars are always
        in sync across AccX, AccY, and AccZ.
  \item Release the mouse button to set the position.
\end{enumerate}

\subsubsection{Method B — Spin Box Entry}

\begin{enumerate}[leftmargin=*, label=\textbf{\arabic*.}]
  \item In the \textbf{Crop Range} panel on the left, find the
        \textbf{Start Time (s)} and \textbf{End Time (s)} spin boxes.
  \item Click a spin box and type a value directly, or use the
        \key{$\uparrow$} / \key{$\downarrow$} arrows to step by 0.1 s.
  \item The trim bars on the plot move instantly to match.
\end{enumerate}

\begin{notebox}
The \textbf{live selection info} below the spin boxes shows the exact
number of records selected and the precise time range in real time as
you drag or type.
\end{notebox}

\subsection{Saving the Cropped File}

\resetsteps
\step{Once the desired range is selected, click \btn{Crop \& Save}.}
\step{A file save dialog opens, pre-filled with a suggested filename of the form:
\begin{lstlisting}
originalname_crop_12.3-45.6s.bin
\end{lstlisting}
Change the filename or destination if needed, then click \key{Save}.}
\step{The console confirms the number of records written and the file path.}
\step{A popup asks: \textit{``Would you like to add it to the FFT Analyzer
      tab?''}  Click \btn{Yes} to load the cropped file into the FFT Analyzer
      immediately, or \btn{No} to skip.}

\begin{warnbox}
Ensure \textbf{Start Time $<$ End Time}; the application will reject the
crop and log an error otherwise.
\end{warnbox}

\subsection{Navigation Toolbar}

The matplotlib toolbar above the plot provides standard navigation:

\begin{center}
\begin{tabular}{ll}
\toprule
\textbf{Button} & \textbf{Action}\\
\midrule
Home        & Reset to the original full view \\
Back / Forward & Step through view history \\
Pan (hand)  & Click-drag to pan the plot \\
Zoom (magnifier) & Draw a rectangle to zoom into that region \\
Save (floppy) & Export the current plot as an image \\
\bottomrule
\end{tabular}
\end{center}

\begin{tipbox}
Zoom and pan are \textbf{independent} of the trim bars.  You can zoom
into a noisy region to precisely position a trim bar without affecting
the crop range.
\end{tipbox}

% ============================================================
\section{FFT Analyzer Tab}
% ============================================================

The FFT Analyzer tab loads one or more \texttt{.bin} files and overlays
their frequency-domain (FFT) and time-domain plots on a shared canvas.
Up to \textbf{10 files} can be loaded simultaneously.

\subsection{Loading Files}

\resetsteps
\step{Click \btn{Add Files} in the \textbf{Files} group.}
\step{A multi-select file browser opens.  Hold \key{Ctrl} or \key{Shift} to
      select multiple \texttt{.bin} files at once, then click \key{Open}.}
\step{Each loaded file appears as a row in the file list, showing:
      a coloured swatch (unique per file), the truncated filename, and a
      visibility checkbox.}
\step{The \textbf{file count} label updates (e.g., ``3 files loaded'').}
\step{Click \btn{Analyze} to draw the plots.}

\subsubsection{Removing All Files}

Click \btn{Clear} to unload all files and reset the canvas.

\subsection{FFT Settings Panel}

\begin{center}
\begin{tabularx}{\linewidth}{lX}
\toprule
\textbf{Control} & \textbf{Description}\\
\midrule
Smoothing Window    & Moving-average window (samples) applied to the FFT
                      magnitude spectrum.  Higher values = smoother curve.
                      Default: 10. \\[4pt]
Auto Y-limit        & When checked (default), the Y-axis upper limit is
                      automatically set to $1.3 \times$ the peak magnitude
                      across all loaded files. \\[4pt]
Y-Axis Limit        & Manual Y-axis ceiling, used only when Auto Y-limit is
                      unchecked.  Default: 0.5 g. \\[4pt]
Log frequency scale & Switches the FFT X-axis to a logarithmic scale,
                      useful for viewing a wide frequency range. \\[4pt]
Frequency Range     & Min and max frequency (Hz) displayed on the FFT
                      X-axis.  Default: 0 -- 50 Hz. \\
\bottomrule
\end{tabularx}
\end{center}

\subsection{Show Plots Panel}

Two checkboxes control which plot types are rendered:

\begin{itemize}
  \item \textbf{FFT (Frequency Domain)} — Shows the single-sided magnitude
        spectrum for each axis.  Y-axis unit: $|$FFT$|$ (g).
  \item \textbf{Time Domain} — Shows raw acceleration versus time.
\end{itemize}

Unchecking one option doubles the width of the remaining plots.
At least one must remain checked.

\subsection{Visibility Checkboxes}

Each file row has a checkbox.  Unchecking a file \textbf{instantly hides}
it from all plots without requiring a re-click of \btn{Analyze}.  The
Y-axis auto-limit recalculates based on the remaining visible files.

\subsection{Maximize / Restore}

Each subplot has a small overlay button in its top-right corner:

\begin{itemize}
  \item \textbf{$\nearrow$ (Maximize)} — Expands that single plot to fill
        the entire canvas for closer inspection.
  \item \textbf{$\swarrow$ (Restore)} — Returns to the full 3-axis grid view.
\end{itemize}

\subsection{Data Cursor}

The data cursor lets you click on any FFT plot to read the exact frequency
and magnitude at a point.

\resetsteps
\step{Click the \textbf{$\oplus$ cursor button} (⌖) in the toolbar to
      enter cursor mode.  The button highlights when active.}
\step{Left-click anywhere on an FFT subplot.  A callout tip snaps to the
      nearest data point on the closest plotted line and shows:
\begin{lstlisting}
f   = 12.345 Hz
|A| = 0.00312 g
\end{lstlisting}}
\step{To dismiss a tip, \textbf{right-click} anywhere inside that subplot.}
\step{Click the cursor button again to exit cursor mode (all tips are cleared).}

\begin{notebox}
The data cursor only works on \textbf{FFT subplots}.  Clicking on
time-domain subplots while in cursor mode has no effect.  The cursor
is also automatically suspended when zoom or pan mode is active.
\end{notebox}

\subsection{FFT Algorithm}

The FFT computation matches the original MATLAB \texttt{VibeDriver.m}
normalisation:

\begin{align}
  \text{mag}[k] &= \left|\frac{X[k]}{N}\right|, \quad k = 0, 1, \ldots, \tfrac{N}{2}\\
  \text{mag}[k] &\times= 2, \quad k = 1, \ldots, \tfrac{N}{2}-1
\end{align}

where $X$ is the DFT of the de-trended signal and $N$ is the signal length.
\texttt{scipy.signal.detrend} removes the DC offset and linear drift before
the FFT.

% ============================================================
\section{Spectrogram Tab}
% ============================================================

The Spectrogram tab generates a high-resolution time-frequency heatmap of
the AccX channel using the \textbf{white-hot} colourmap (dark navy → blue →
orange → red → white), matching the MATLAB \texttt{Spectrogram\_SFFT.m}
output.

\subsection{Loading a File}

\resetsteps
\step{Click \btn{Load BIN File}.}
\step{Select a \texttt{.bin} file.  The left panel shows the filename,
      record count, time range, and sample rate.}
\step{The \btn{Analyze} button becomes active.}

\subsection{Spectrogram Settings}

\begin{center}
\begin{tabularx}{\linewidth}{lX}
\toprule
\textbf{Control} & \textbf{Description}\\
\midrule
Window Size      & Length of each FFT window (samples).  Larger windows give
                   better \textbf{frequency resolution} but worse time
                   resolution.  Range: 64--16384.  Default: 1024. \\[4pt]
Magnitude Limit (g) & The upper bound of the colour scale.  Data points at or
                   above this value are rendered \textbf{white} (hottest
                   colour).  Increase if most of the plot looks white;
                   decrease if it looks mostly dark.  Default: 1.2 g. \\[4pt]
Freq Range (Hz)  & The Y-axis frequency range displayed.  Only affects the
                   visible area — all frequencies are still computed.
                   Default: 0 -- 40 Hz. \\
\bottomrule
\end{tabularx}
\end{center}

\begin{notebox}
\textbf{Overlap} and \textbf{NFFT} are computed automatically:
\begin{itemize}
  \item Overlap $= \lfloor\text{Window} / 2\rfloor$ \;(50\,\% overlap)
  \item NFFT $= \max(4096,\; \text{Window} \times 4)$
\end{itemize}
These defaults provide good time-frequency balance for typical
flight/vibration recordings without requiring manual tuning.
\end{notebox}

\subsection{Running the Analysis}

\resetsteps
\step{Adjust the settings if needed (or leave defaults).}
\step{Click \btn{Analyze}.}
\step{The canvas shows \textit{``Computing spectrogram\ldots''} and the buttons
      are temporarily disabled.  For long recordings this may take a few
      seconds — the application remains \textbf{fully responsive} during
      computation.}
\step{When done, the white-hot spectrogram appears:
      \begin{itemize}
        \item X-axis: Flight time (s)
        \item Y-axis: Frequency (Hz)
        \item Colour scale: Magnitude (g) from dark (low) to white (high)
      \end{itemize}}

\subsection{Reading the Spectrogram}

\begin{itemize}
  \item \textbf{Bright / white regions} indicate high vibration energy at
        that frequency and time — potential resonance or structural response.
  \item \textbf{Dark / navy regions} indicate low energy — quiet periods or
        frequencies with little activity.
  \item \textbf{Horizontal bands} (constant frequency, persistent) suggest
        structural resonance modes.
  \item \textbf{Vertical bright stripes} (across all frequencies, brief time)
        suggest short impulse events.
\end{itemize}

\subsection{Adjusting the Colour Scale}

If the plot appears \textbf{mostly white}: reduce the \textbf{Magnitude Limit}
and re-run.  If the plot appears \textbf{mostly dark}: increase it and re-run.
A good starting point is just above the background noise floor.

% ============================================================
\section{Keyboard \& Mouse Reference}
% ============================================================

\begin{center}
\begin{tabularx}{\linewidth}{llX}
\toprule
\textbf{Location} & \textbf{Input} & \textbf{Action}\\
\midrule
Password dialog   & \key{Enter}        & Submit password \\
Spin boxes        & \key{$\uparrow$} / \key{$\downarrow$} & Increment / decrement value \\
Any plot          & Scroll wheel       & Zoom in/out (when zoom mode active) \\
Crop plot         & Left-click + drag near trim bar & Move start or end bar \\
Crop plot         & Mouse hover near bar & Cursor changes to resize arrow \\
FFT plot (cursor mode) & Left-click   & Place data-cursor tip \\
FFT plot (cursor mode) & Right-click  & Dismiss tip on that subplot \\
Toolbar           & Click Zoom, then drag & Draw zoom rectangle \\
Toolbar           & Click Pan, then drag  & Pan the view \\
Toolbar           & Home icon             & Reset zoom/pan to full view \\
\bottomrule
\end{tabularx}
\end{center}

% ============================================================
\section{Troubleshooting}
% ============================================================

\begin{center}
\begin{tabularx}{\linewidth}{lX}
\toprule
\textbf{Symptom} & \textbf{Solution}\\
\midrule
\texttt{vibe} not recognised & Contact \texttt{Akshat\_EA} to restore the
                               PowerShell alias. \\[4pt]
Application does not launch  & Ensure Python 3 is installed and
                               \texttt{vibe} points to the correct path.
                               Try running \texttt{python main.py} manually
                               from the \texttt{VibrationAnalysis} folder. \\[4pt]
\textit{``File has no complete records''} error & The selected file is empty
                               or corrupted.  Verify the file size is a
                               multiple of 18 bytes. \\[4pt]
Spectrogram takes very long  & Reduce the \textbf{Window Size} (e.g., to 512)
                               for a faster, lower-resolution result. \\[4pt]
FFT plot looks flat / all zero & The file may contain NaN or corrupt records.
                               The application substitutes zeros for corrupt
                               values automatically. \\[4pt]
Three wrong password attempts & Re-launch with \texttt{vibe} and contact
                               \texttt{Akshat\_EA} for the correct password. \\[4pt]
Theme colours not updating   & This is a known cosmetic issue on some
                               display drivers.  Toggle the theme twice
                               to force a repaint. \\
\bottomrule
\end{tabularx}
\end{center}

% ============================================================
\section{Contact \& Support}
% ============================================================

For access credentials, technical issues, or feature requests, contact:

\begin{center}
\Large\textbf{\color{accent} Akshat\_EA}
\end{center}

\vspace{0.5cm}

\begin{notebox}
This application and its documentation are developed and maintained
exclusively by \textbf{Akshat\_EA}.  Unauthorised distribution or
modification is not permitted.
\end{notebox}

% ── Back cover ───────────────────────────────────────────────
\newpage
\pagecolor{headerbg}
\color{headertext}
\vspace*{\fill}
\begin{center}
  {\Large\bfseries Vibration Analysis Suite}\\[0.4cm]
  {\color{muted} User Manual — Version 1.0}\\[1.5cm]
  {\color{muted}\small Developed by\;}\textbf{\color{accent}\large Akshat\_EA}\\[0.4cm]
  {\color{muted}\small \today}
\end{center}
\vspace*{\fill}

\end{document}
